\section{Gestione degli errori}

\subsubsection{Errori ed eccezioni}

Tutte le computazioni, come le funzioni chiamate, possono fallire prematuramente.
A volte il risultato attenon non può essere genreato (log di 3), a volte non possono
generare gli effeti collaterare( non cè scritto nulla sul file).
Il paggio: ha iniziato a fare degli effetti collatewrali ma non li fa
tutti (eg salvo metaà immagine poi il disco full).

Questi fallimenti possono essere dovuti a cause attese,
come un numero al posto di una stringa.

Altri errori sono legati al sistema di esecuzione soggiacente (il pc, il cloud)
in realtà le capacità sono limitate.

Lo spazio finisce sia SSD sia RAM. A volte si rompono fisicamente.
A volte ho perso la connessione di rete.

Indipendentemente dalle cause che li hanno prodotti, i fallimenti possono essere
catalogati in due gruppi principali

Recuperabili quelli che non hanno comprosmeesso lo stato del programma,
nel senso che possiamo tornare indietro, ho riportato lo stato incositene
ad uno stato consistente

Altri sono recuperabili, quindi ipediscono qualuncuq operazione di
recupero. Ed occorre interrompere il processi

Solitamente voglimao uscire dal processo in maniera pulita, negli
altri occorre mettere in atto una strategia di ripristio di stato.

Il punto in cui rilevo che c'è qualcosa che non va
può essere difficile trovarlo, magari in una fuznion (eg se il
diso è piano è la sys.write che fallisce ma non sa cosa farsene)
Occorre che quella sys call li al suo chiamante che non è andata a buon
fine.


Ci serve che quando rilevo un fallimento possa interrompere la compiutazion
corrente e ritornare al chiamanerte comunicando cosa c'è stato di andato storto.

Ma questo aumenta la complessità ciclomatica del codice!



\subsubsection{Spazio di scelta}
Dato un tipo di errore che posso farci io, che spazio di scelt ho?

In alcuni casi può essere causato da una causa treansitoria (vedi locked di un file, aspetto un po e top)
ooopures server di rete non raggiungibile ora, dopo quanto tempo posso
riporvare (exponential backoff per al rete).

Altra domanda, se non ha senso risporvare l'esecuzzione, posso
continuare comunque con altre operazioni?
Eg leggo un file se una file è rotta la salto

se devo leggere il file di config, non lo trovo, uso dei default.

Si rompe qualcosa che deriva dall'utente, llo posso convolgere?

Posso isolare l'errore e mantere consistente tutto il resto, vedi il DB. faccio
il rollback di una sola transazione, non di tutto il DB.

Si è verificato un Bug, allora la pianto subito.


\subsection{Supporto sintattico alla gestione degli errori}

Nei linguaggi C, java, python, c++  viene introditta il concetto di
eccezione, tipo di dato che descrive un errore che viene lasnciato con throw
e viene restituito al chimante.
il chiamnte deve averre incluso la chiamta in un blocco try catch
altrimenti risale fino al main o al primo try catch che lo gestisce
altirmenti il pgm muore.

QUesta idea che sembrava fantasmagoriica ma ha dato solo problemi,
la fuznione va fino in fondo o no?

Trows some exception se in java un metodo può lanciare un eccezione.

error --> rompe la JVM eg stack overflow e sono unrecoverable (ti trapano lo stesso il blocco try catvh)
Exception --> in 2 famiglie le run time  -> non catchabili nullpointer ececetopn alloa non servono
e tutte le altre

in più creo probli nello stack, devo tenere tracce del contesto delle eccezioni.
In liinux non è possibile scrivere codice del kernel in C++ in linux(?), in C c'è un codice errore, te lo scelgi te.

In Rust descrivo un tipo che si chiama Resiult

ok() che racchiude il risultato
err() che racchiude l'errore

è una enum generica Enum Result<T, E>.


la macro panic va usata con cautela.

\begin{minted}[breaklines]{rust}
enum Result<T, E> {
    Ok(T),
    Err(E),
  }
\end{minted}

\subsection{I limiti della gestione delle eccezioni}



\subsection{Gestioni delle eccezioni in Rust}



\begin{minted}[breaklines]{rust}
fn read_file(name: &str)-> Result<String,io::Error> {
    let r1 = File::open(name);
    let mut file = match r1 {
        Err(why) => return Err(why),
        Ok(file) => file,
      };
    let mut s = String::new();
    let r2 = file.read_to_string(&mut s);
    match (r2) {
        Err(why) => Err(why),
        Ok(_) => Ok(s),
      }
  }
\end{minted}

\subsection{Elaborare i risultati}

Il metdoo ok trasforma il resutl in una Option,
err trasforma il result in una Option


il map trasorma il valroe buono e trasforma in un valore
in un nuovo l'errore. ritorna un Resutl

Contains permette di verficare se cntien un valore == è booleana.

unwrap prova a dire per me è ok, ti ritonra l'errore altriemtni panic,
non va benen perdiamo il controlllo

expect, al posto du usare il messaggio dei errore che vi è nel messsaggio,
lo decisiamo noi.

\subsection{Gestire gli errori}

\subsection{Panic}


\subsection{Ignorare gli errori}


\subsection{Propagare gli errori}
nelle fuznion che scriviammo non abbiamo contesto sufficente per gestire

l'errore quindi il ritorno al chiamnre è la più semplice.

rust mette a disposizione l'operatore ? che applicato a RESULT

se ritorna ok, da il valroe altriemnti, il ? comporta un
ritorno immediato dell'errore al chiamante il quale la deve gestire.


\begin{minted}[breaklines]{rust}
fn read_file(name: &str) -> Result<String,io::Error> {
  let mut file = File::open(name)?;
  let mut s = String::new();
  file.read_to_string(&mut s)?;
  Ok(s)
}


\end{minted}

\subsection{Altri modi di esprimere il fallimento}
Posso convertire l'error in un Otpion

\subsection{Propagare errori eterogenei}
il ? è comodo ma ha un vincolo.

1. se la funziona deve dire che ritorna un Result di che cosa?
ma gli errori possono essere eterogenei nella stessa funizone. Eg
data una stringa voglio convertire in un numero
2 errori 1. il disco IO e 2. pars error. Non posso fare un result
che ritorna a caso ma solo 1 tipo di errore


Come gestiamo, usimo il polimorfismo, diciamo che la fuznione
nella result c'è il tipo che ci interessa se è ok altrimenti un Box
che implementa il trattto error box<dyn error::Error> e può essere costrto
da qualsiasi tipo di errore.
Pronlema se il chiamante vuole gestire in modo diverso i due errori, non può farlo
si dovrebbe usare un downcastrest

\begin{minted}[breaklines]{rust}
fn sum_file(path: &Path) -> Result<i32, Box<dyn error::Error>> {
  let mut file = File::open(path) ? ; // io::Error -> Box<dyn error::Error>
  let mut contents = String::new();
  file.read_to_string(&mut contents) ? ; // io::Error -> Box<dyn error::Error>
  let mut sum = 0;
  for line in contents.lines() {
    sum += line.parse::<i32>() ? ; // ParseIntError -> Box<dyn error::Error>
  }
  Ok(sum)
}

fn handle_sum_file_errors(path: &Path) {
  match sum_file(path) {
    Ok(sum) => println!("sum is {}", sum),
    Err(err) => {
      if let Some(e) = err.downcast_ref::<io::Error>() {...} //tratto io::Error
      else if let Some(e) = err.downcast_ref::<ParseIntError>() {...} //tratto ParseIntError
      else { unreachable!(); } //non può capitare
    }
  }
}
\end{minted}


\subsection{Propagare errori eterogenei}

due strede con Enum, creo un

problem, se può essere accettabile con 


thiserror per le librerie ed anyhow per i programmi applicativi.
